\chapter{Przebieg zmienności funkcji}
\label{ma1:ch:function-variation}

Badanie przebiegu zmienności polega na powiązaniu informacji o dziedzinie, granicach, znakach pochodnych i ekstremach z geometrycznym obrazem funkcji. Sama tabela znaków pochodnej nie rozstrzyga jeszcze o globalnym zachowaniu funkcji, zwłaszcza gdy jej dziedzina nie jest przedziałem.

\section{Punkty stacjonarne i kryterium pierwszej pochodnej}
\begin{definicja}
\label{ma1:def:stationary-critical}
Niech $f:D\to\R$, gdzie $D\subseteq\R$. Punkt wewnętrzny $a$ zbioru $D$ nazywamy stacjonarnym, jeśli $f'(a)$ istnieje i $f'(a)=0$. Punktem krytycznym nazywamy tutaj punkt wewnętrzny dziedziny, w którym pochodna jest zerowa albo nie istnieje. Końce dziedziny rozpatrujemy oddzielnie.
\end{definicja}
Z twierdzenia Fermata \ref{ma1:thm:fermat} każde ekstremum lokalne w punkcie wewnętrznym dziedziny musi być punktem krytycznym: jeśli pochodna istnieje, to jest zerowa, a jeśli nie istnieje, punkt należy do drugiej klasy z definicji. Implikacja odwrotna jest fałszywa. Funkcja $x\mapsto x^3$ ma pochodną $3x^2$, zatem punkt $0$ jest stacjonarny, ale nie jest ekstremum, ponieważ dla każdego $h>0$ zachodzi $(-h)^3<0<h^3$.

\begin{twierdzenie}[Kryterium znaku pierwszej pochodnej]
\label{ma1:thm:first-derivative-test}
Niech $f$ będzie ciągła w punkcie $a$ i różniczkowalna na $(a-\delta,a)$ oraz $(a,a+\delta)$ dla pewnego $\delta>0$. Jeśli $f'<0$ po lewej stronie $a$ i $f'>0$ po prawej, to $a$ jest punktem ścisłego minimum lokalnego. Przy odwrotnych znakach otrzymujemy ścisłe maksimum lokalne. Jeżeli $f'>0$ na obu przedziałach albo $f'<0$ na obu przedziałach, to $a$ nie jest punktem ekstremum lokalnego.
\end{twierdzenie}
\begin{proof}
Dla $a-\delta<x<a$ funkcja jest ciągła na $[x,a]$ (w punkcie $a$ z założenia, w pozostałych punktach z różniczkowalności) i różniczkowalna na $(x,a)$. Twierdzenie Lagrange'a \ref{ma1:thm:lagrange} daje $c\in(x,a)$ takie, że
\[
 f(a)-f(x)=f'(c)(a-x).
\]
Ponieważ $a-x>0$, znak różnicy jest znakiem pochodnej po lewej. Dla $a<x<a+\delta$ te same założenia na $[a,x]$ dają $d\in(a,x)$ i
$f(x)-f(a)=f'(d)(x-a)$. Zatem przy znakach $-,+$ obie wartości $f(x)$ są większe od $f(a)$, a przy znakach $+,-$ obie są mniejsze. Przy znakach $+,+$ wartości po lewej są mniejsze, a po prawej większe od $f(a)$; przy $-,-$ jest odwrotnie. To dowodzi wszystkich czterech tez.
\end{proof}

\begin{przyklad}
\label{ma1:ex:critical-nondifferentiable}
Funkcja $f(x)=|x|$ jest ciągła w zerze, a jej pochodna wynosi $-1$ dla $x<0$ i $1$ dla $x>0$. Ilorazy różnicowe w zerze wynoszą odpowiednio $-1$ i $1$, więc pochodna w zerze nie istnieje. Kryterium \ref{ma1:thm:first-derivative-test} daje ścisłe minimum. Sama ciągłość w punkcie krytycznym jest istotna: funkcja $g(x)=|x|$ dla $x\ne0$, $g(0)=1$, ma takie same pochodne na obu półosiach, ale dla $0<|x|<1$ zachodzi $g(x)<g(0)$, więc w zerze ma ścisłe maksimum, nie minimum.
\end{przyklad}

\section{Kryteria wykorzystujące wyższe pochodne}
\begin{twierdzenie}
\label{ma1:thm:higher-derivative-extrema}
Niech $f$ będzie określona na otwartym przedziale zawierającym $a$ i niech $m\ge1$ będzie liczbą całkowitą. Załóżmy, że $f$ jest $(m-1)$ razy różniczkowalna w otoczeniu $a$, a $f^{(m-1)}$ jest różniczkowalna w $a$, przy czym
\[
 f^{(j)}(a)=0\quad(1\le j<m),\qquad f^{(m)}(a)\ne0.
\]
Jeżeli $m$ jest parzyste, to w $a$ występuje ścisłe minimum lokalne dla $f^{(m)}(a)>0$ i ścisłe maksimum lokalne dla $f^{(m)}(a)<0$. Jeżeli $m$ jest nieparzyste, to w $a$ nie ma ekstremum lokalnego.
\end{twierdzenie}
\begin{proof}
Twierdzenie Taylora z resztą Peana \ref{ma1:thm:taylor-peano} można zastosować na mocy założeń o istnieniu pochodnych. Po zniknięciu współczynników niższych rzędów dostajemy dla $h\ne0$ dostatecznie małego
\[
 f(a+h)-f(a)=h^m(c+\varepsilon(h)),\qquad
 c=\frac{f^{(m)}(a)}{m!},\qquad \varepsilon(h)\longrightarrow0.
\]
Ponieważ $c\ne0$, istnieje $\eta>0$ taka, że dla $0<|h|<\eta$ mamy $|\varepsilon(h)|<|c|/2$. W szczególności $c+\varepsilon(h)$ ma znak $c$. Dla parzystego $m$ czynnik $h^m$ jest dodatni niezależnie od strony, więc różnica ma zawsze znak $c$. Dla nieparzystego $m$ czynnik $h^m$ ma znak $h$, zatem różnica przyjmuje wartości obu znaków dowolnie blisko $a$. Otrzymujemy żądane klasyfikacje.
\end{proof}
W szczególności, jeśli $f'(a)=0$ i istnieje $f''(a)>0$, mamy ścisłe minimum; jeśli $f''(a)<0$, mamy ścisłe maksimum. Warunek $f''(a)=0$ niczego nie rozstrzyga. Funkcje $x^4$, $-x^4$ i $x^3$ mają w zerze pierwszą i drugą pochodną zerową, lecz odpowiednio ścisłe minimum, ścisłe maksimum i brak ekstremum, co wynika z twierdzenia \ref{ma1:thm:higher-derivative-extrema} dla rzędów $4,4,3$.

\begin{przyklad}
\label{ma1:ex:flat-extrema}
Nie zawsze istnieje pierwszy niezerowy współczynnik wyższego rzędu. Niech $H$ będzie funkcją z przykładu \ref{ma1:ex:taylor-flat}, czyli $H(0)=0$ i $H(x)=e^{-1/x^2}$ dla $x\ne0$. Wszystkie pochodne dodatnich rzędów funkcji $H$, $-H$ oraz $xH(x)$ w zerze są równe zeru. Dla ostatniej funkcji wynika to ze wzoru Leibniza \ref{ma1:thm:leibniz}: dla $n\ge1$ jej pochodna w zerze jest równa $0\cdot H^{(n)}(0)+nH^{(n-1)}(0)=0$. Jednak $H(x)>0$ poza zerem, więc $H$ ma ścisłe minimum, $-H$ ścisłe maksimum, a $xH(x)$ ma znak $x$ i nie ma ekstremum w zerze. Same wartości wszystkich pochodnych w punkcie nie muszą zatem wyznaczać typu ekstremum.
\end{przyklad}

\section{Asymptoty}
\begin{definicja}
\label{ma1:def:asymptotes}
Prostą $x=a$ nazywamy asymptotą pionową wykresu $f$, jeśli co najmniej jedna z granic jednostronnych przy $x\to a$ w dziedzinie funkcji wynosi $+\infty$ albo $-\infty$ (rozpatrujemy tylko strony, z których $a$ jest punktem skupienia dziedziny).
Jeśli dziedzina jest nieograniczona z prawej albo z lewej, prostą $y=mx+b$ nazywamy asymptotą w odpowiednim kierunku, gdy
\[
 \lim_{x\to+\infty}\bigl(f(x)-mx-b\bigr)=0
 \quad\text{lub}\quad
 \lim_{x\to-\infty}\bigl(f(x)-mx-b\bigr)=0.
\]
Dla $m=0$ jest to asymptota pozioma, a dla $m\ne0$ ukośna. Granice rozumiemy względnie do dziedziny.
\end{definicja}
\begin{stwierdzenie}
\label{ma1:prop:asymptote-coefficients}
W każdym z dwóch kierunków nieskończoności, w którym dziedzina jest nieograniczona, istnieje asymptota postaci $y=mx+b$ wtedy i tylko wtedy, gdy istnieją skończone granice
\[
 m=\lim_{x\to\pm\infty}\frac{f(x)}x,\qquad
 b=\lim_{x\to\pm\infty}(f(x)-mx).
\]
W obu granicach wybieramy ten sam kierunek. Asymptota w danym kierunku jest wówczas jedyna.
\end{stwierdzenie}
\begin{proof}
Jeżeli $f(x)=mx+b+r(x)$, gdzie $r(x)\to0$ w ustalonym kierunku, to dla $x\ne0$ mamy $f(x)/x=m+b/x+r(x)/x$. Oba ostatnie składniki dążą do zera, drugi dlatego, że $r(x)$ jest ograniczona dla dostatecznie dużego $|x|$. Ponadto $f(x)-mx=b+r(x)\to b$. Odwrotnie, jeżeli istnieją podane skończone granice, to $f(x)-mx-b\to0$, co jest definicją asymptoty. Jednoznaczność wynika z jednoznaczności obu granic: każda asymptota musi mieć właśnie tak wyznaczone współczynniki.
\end{proof}
\begin{przyklad}
\label{ma1:ex:asymptote-slope-only}
Dla $f(x)=x+\sqrt{x}$ na $(0,\infty)$ mamy $f(x)/x=1+1/\sqrt{x}\to1$ przy $x\to+\infty$. Jednak $f(x)-x=\sqrt{x}\to+\infty$, więc druga granica nie jest skończona i nie ma asymptoty ukośnej ani poziomej. Istnienie pierwszej granicy nie wystarcza.
\end{przyklad}

\section{Schemat badania funkcji i ekstrema globalne}
Najpierw wyznaczamy dziedzinę, miejsca zerowe, znak funkcji oraz ewentualną symetrię. Następnie badamy ciągłość i granice na krańcach składowych dziedziny. Wyznaczamy pochodną, punkty krytyczne i przedziały, na których możemy zastosować twierdzenie o monotoniczności \ref{ma1:thm:derivative-monotonicity}. Klasyfikujemy ekstrema, sprawdzamy końce dziedziny oraz granice w nieskończoności, wyznaczamy asymptoty i zbiór wartości. Dopiero na tej podstawie szkicujemy wykres. Przy skomplikowanej dziedzinie lub nieskończonym zbiorze punktów krytycznych skończona tabela znaków nie musi być możliwa; schemat nie jest algorytmem gwarantującym skończone rozwiązanie dla dowolnej funkcji. Szczegółowe kryteria wypukłości i punkty przegięcia omówimy w kolejnym rozdziale.

\begin{stwierdzenie}
\label{ma1:prop:global-extrema-candidates}
Niech $f:[a,b]\to\R$ będzie ciągła, gdzie $a<b$. Wszystkie punkty osiągania minimum lub maksimum globalnego należą do zbioru złożonego z końców $a,b$ oraz punktów krytycznych w $(a,b)$. Jeżeli zbiór ten jest skończony, wystarczy porównać wartości funkcji w jego punktach.
\end{stwierdzenie}
\begin{proof}
Twierdzenie o osiąganiu kresów \ref{ma1:cor:interval-extrema} zapewnia istnienie punktów osiągania minimum i maksimum, ponieważ $f$ jest ciągła na przedziale domkniętym i ograniczonym. Jeśli taki punkt leży we wnętrzu, jest także ekstremum lokalnym, więc na mocy twierdzenia Fermata \ref{ma1:thm:fermat} pochodna jest tam zerowa, o ile istnieje. W pozostałym przypadku punkt jest krytyczny z definicji; końce są jawnie uwzględnione. Porównanie wartości na skończonym zbiorze zawierającym wszystkie punkty osiągania kresów daje te same minimum i maksimum co na całej dziedzinie.
\end{proof}
Na przedziale otwartym lub nieograniczonym kres nie musi być osiągany: funkcja $x\mapsto x$ na $(0,1)$ ma kres dolny $0$ i górny $1$, lecz żadna z tych liczb nie należy do jej zbioru wartości $(0,1)$.

\section{Przykłady pełnego badania}
\begin{przyklad}
\label{ma1:ex:variation-rational}
Rozważmy $f(x)=x+1/x$ na $D=\R\setminus\{0\}$. Funkcja jest nieparzysta, gdyż $f(-x)=-f(x)$. Nie ma miejsc zerowych: $f(x)=(x^2+1)/x$, a licznik jest zawsze dodatni, więc znak funkcji jest znakiem $x$. Funkcja jest ciągła i różniczkowalna na obu składowych dziedziny. Z granic składników otrzymujemy
\[
 \lim_{x\to0^-}f(x)=-\infty,\qquad
 \lim_{x\to0^+}f(x)=+\infty.
\]
Prosta $x=0$ jest asymptotą pionową, a $y=x$ asymptotą ukośną w obu kierunkach, ponieważ $f(x)-x=1/x\to0$ dla $x\to\pm\infty$. W szczególności granice $f$ w nieskończoności mają znak odpowiedniej nieskończoności.

Pochodna wynosi
\[
 f'(x)=1-\frac1{x^2}=\frac{(x-1)(x+1)}{x^2}.
\]
Mianownik jest dodatni w dziedzinie. Licznik jest dodatni dla $|x|>1$, ujemny dla $0<|x|<1$ i zerowy dla $x=\pm1$. Z twierdzenia \ref{ma1:thm:derivative-monotonicity} dostajemy tabelę:
\[
\begin{array}{c|ccccccc}
 x&(-\infty,-1)&-1&(-1,0)&0&(0,1)&1&(1,\infty)\\ \hline
 f'&+&0&-&\text{---}&-&0&+\\
 f&\nearrow&-2&\searrow&\text{---}&\searrow&2&\nearrow
\end{array}
\]
W punkcie $-1$ jest ścisłe maksimum lokalne, a w $1$ ścisłe minimum lokalne, na mocy twierdzenia \ref{ma1:thm:first-derivative-test}. Są to również odpowiednio maksimum na ujemnej i minimum na dodatniej półosi. Dla $x>0$ zachodzi $f(x)-2=(x-1)^2/x\ge0$. Każda wartość $y>2$ jest osiągana na $(1,\infty)$: wybieramy $B>1$ z $f(B)>y$, co umożliwia granica w $+\infty$, i stosujemy twierdzenie o wartości pośredniej \ref{ma1:thm:intermediate-value} do ciągłej funkcji na $[1,B]$. Wartość $2$ jest osiągana w $1$. Nieparzystość daje zbiór wartości na ujemnej półosi. Ostatecznie
\[
 f(D)=(-\infty,-2]\cup[2,\infty).
\]
Na całej dziedzinie nie ma ani maksimum, ani minimum globalnego, ponieważ funkcja jest nieograniczona z obu stron.
\end{przyklad}
\begin{center}
\begin{tikzpicture}[x=1.1cm,y=.85cm,>=Stealth]
 \draw[->] (-4.4,0)--(4.4,0) node[right] {$x$};
 \draw[->] (0,-5.4)--(0,5.4) node[above] {$y$};
 \draw[gray,dashed] (-4.2,-4.2)--(4.2,4.2)
 node[above left] {$y=x$};
 \draw[blue!65!black,thick,domain=.22:4,samples=140] plot (\x,{\x+1/\x});
 \draw[blue!65!black,thick,domain=-4:-.22,samples=140] plot (\x,{\x+1/\x});
 \fill (1,2) circle (2pt) node[above right] {$(1,2)$};
 \fill (-1,-2) circle (2pt) node[below left] {$(-1,-2)$};
 \node[blue!65!black,above left] at (2.6,3.0) {$f$};
 \node[gray,right] at (0,-4.9) {$x=0$};
\end{tikzpicture}

\small Rysunek: $f(x)=x+1/x$; gałęzie rozdziela asymptota pionowa, a asymptota ukośna jest wspólna dla obu kierunków.
\end{center}

\begin{przyklad}
\label{ma1:ex:variation-exp}
Niech $g(x)=x^2e^{-x}$ na $\R$. Funkcja jest gładka na całej prostej, nieujemna, a jej jedyne miejsce zerowe to $0$, ponieważ $e^{-x}>0$. Nie jest ani parzysta, ani nieparzysta: $g(1)=e^{-1}$ oraz $g(-1)=e$, a $e>1$. Przy $x\to-\infty$ mamy $g(x)\ge x^2$ dla $x<0$, więc $g(x)\to+\infty$. Dla $x>0$ ze wzoru Taylora z resztą Lagrange'a \ref{ma1:thm:taylor-remainders} dla funkcji wykładniczej w rzędzie $3$ mamy $e^x\ge x^3/3!$, gdyż wszystkie składniki i reszta są dodatnie. Stąd $0\le g(x)\le 6/x\to0$ przy $x\to+\infty$. Prosta $y=0$ jest asymptotą poziomą w tym kierunku; w kierunku $-\infty$ nie ma asymptoty postaci $y=mx+b$, bo dla $x<0$ mamy $g(x)/x=xe^{-x}\le x\to-\infty$. Asymptot pionowych nie ma z ciągłości na $\R$.

Z reguł iloczynu i łańcuchowej wynika
\[
 g'(x)=2xe^{-x}-x^2e^{-x}=x(2-x)e^{-x}.
\]
Ponieważ czynnik wykładniczy jest dodatni, znaki pochodnej na kolejnych przedziałach $( -\infty,0)$, $(0,2)$ i $(2,\infty)$ to $-,+,-$. Twierdzenia \ref{ma1:thm:derivative-monotonicity} i \ref{ma1:thm:first-derivative-test} dają ścisłe zmniejszanie, zwiększanie i zmniejszanie oraz ścisłe minimum w $0$ i ścisłe maksimum lokalne w $2$ o wartości $4e^{-2}$. Minimum w $0$ jest globalne z nieujemności funkcji, maksimum w $2$ nie jest globalne, gdyż $g$ jest nieograniczona z góry. Zbiór wartości to $[0,\infty)$: dla $y>0$ wybieramy $B<0$ z $g(B)>y$ i stosujemy twierdzenie o wartości pośredniej na $[B,0]$, gdzie funkcja jest ciągła.
\end{przyklad}
\begin{center}
\begin{tikzpicture}[x=1.55cm,y=3.2cm,>=Stealth]
 \draw[->] (-1.15,0)--(5.4,0) node[right] {$x$};
 \draw[->] (0,-.05)--(0,1.32) node[above] {$y$};
 \draw[blue!65!black,thick,domain=-.72:5.1,samples=140] plot (\x,{\x*\x*exp(-\x)});
 \draw[gray,dotted] (2,0)--(2,{4*exp(-2)});
 \fill (0,0) circle (2pt) node[below left] {$(0,0)$};
 \fill (2,{4*exp(-2)}) circle (2pt) node[above] {$(2,4e^{-2})$};
 \node[below] at (2,0) {$2$};
 \node[blue!65!black,above] at (3.8,.32) {$g$};
\end{tikzpicture}

\small Rysunek: $g(x)=x^2e^{-x}$. Na rysunku widać minimum globalne, maksimum lokalne i dążenie prawej gałęzi do asymptoty poziomej.
\end{center}
